\chapter{Combinatorics}
\label{ch:combinatorics}

\medskip

\textbf{Notation.} $[n]=\{1,\dots,n\}$; $\binom{[n]}{k}$ --- $k$-subsets of $[n]$.
$2^{[n]}$ --- power set; $\binom{n}{k}$ --- binomial coefficient.
$|X|$ --- cardinality; $\mathbb{E}[X]$ --- expectation; $\mathrm{Var}(X)$ --- variance.
$\mathcal{F}$ --- family of sets.
$A+B=\{a+b:a\in A,b\in B\}$ --- sumset; $nA=A+\cdots+A$ ($n$ times).

\medskip

\section{Pigeonhole principle and extremal principle}
\label{sec:pigeonhole}

\subsection{Strengthened pigeonhole principle}
\label{subsec:strengthened-pigeonhole}

$n$ objects into $k$ boxes $\Rightarrow$ some box contains $\ge\lceil n/k\rceil$ objects. Averaging: $\max a_i\ge S/n$, $\min a_i\le S/n$.

\subsection{Extremal principle}
\label{subsec:extremal-principle}

On a finite set of configurations there exist minimizing and maximizing a given function. Choose a configuration with minimal/maximal value in the ``bad'' class, locally vary $\to$ contradiction.

\subsection{Erd\H{o}s--Szekeres theorem}
\label{subsec:erdos-szekeres}

A sequence of $N\ge (r-1)(s-1)+1$ distinct numbers contains an increasing subsequence of length $r$ or a decreasing subsequence of length $s$. The bound is sharp.

\subsection{Dilworth and Mirsky theorems}
\label{subsec:dilworth-mirsky}

A \textit{poset} is a set with a partial order $\le$. A \textit{chain} is a totally ordered subset; an \textit{antichain} is a set of pairwise incomparable elements. In a finite poset: the minimum number of chains covering $P$ equals the width $w(P)$ (Dilworth). The minimum number of antichains covering $P$ equals the height $h(P)$ (Mirsky). Consequence: $|P|\ge ab+1$ $\Rightarrow$ chain of length $a+1$ or antichain of size $b+1$.

\subsection{Erd\H{o}s--Ko--Rado theorem}
\label{subsec:erdos-ko-rado}

$\mathcal{F}\subseteq\binom{[n]}{k}$ is an intersecting family, $n\ge 2k$ $\Rightarrow$ $|\mathcal{F}|\le\binom{n-1}{k-1}$. Equality holds for a star (all $k$-subsets containing a fixed element).

\subsection{Ramsey theorem (binomial bound)}
\label{subsec:ramsey}

$R(s,t)\le R(s-1,t)+R(s,t-1)$, hence $R(s,t)\le\binom{s+t-2}{s-1}$ and $R(s,s)\le 4^{s-1}$.

\medskip

\section{Probabilistic method}
\label{sec:probabilistic-method}

\subsection{Linearity of expectation}
\label{subsec:linearity-expectation}

$\mathbb{E}[\sum X_i]=\sum\mathbb{E}[X_i]$ for any (dependent) integrable $X_i$.

\subsection{First moment method}
\label{subsec:first-moment}

$X\ge 0$, $\mathbb{E}[X]<1$ $\Rightarrow$ $P(X=0)>0$ (existence).

\subsection{Alteration method}
\label{subsec:alteration}

A random construction $C$ is almost good. Remove defects. Obtain a pure construction of size $\ge\mathbb{E}[|C|]-\mathbb{E}[Y]$, where $Y$ is the number of defects.

\subsection{Lov\'asz Local Lemma (LLL)}
\label{subsec:lll}

Events $A_i$, $P(A_i)\le p$, each depends on at most $d$ others, $ep(d+1)\le 1$ $\Rightarrow$ $P(\bigcap\overline{A_i})>0$.

\subsection{Second moment method}
\label{subsec:second-moment}

$P(X=0)\le\frac{\mathrm{Var}(X)}{(\mathbb{E}[X])^2}$. If $\mathrm{Var}(X)=o((\mathbb{E}[X])^2)$, then $X>0$ with probability $\to 1$.

\subsection{Random lower bound for Ramsey numbers}
\label{subsec:ramsey-random}

$R(k,k) > 2^{k/2}$ for $k\ge 3$. Improved by alteration: $\frac{k}{e\sqrt{2}}2^{k/2}$.

\subsection{Chernoff--Hoeffding inequality}
\label{subsec:chernoff}

$X_i\in[a_i,b_i]$ independent, $S=\sum X_i$: $P(|S-\mathbb{E}[S]|\ge t)\le 2\exp(-2t^2/\sum(b_i-a_i)^2)$.

\medskip

\section{Bijective arguments and double counting}
\label{sec:double-counting}

\subsection{Double counting (incidence matrix)}
\label{subsec:double-counting}

For $R\subseteq X\times Y$: $|R|=\sum_x|\{y\}|=\sum_y|\{x\}|$. Consequences: $\sum\deg v=2|E|$, Vandermonde's identity, hockey-stick identity.

\subsection{LYM inequality and Sperner's theorem}
\label{subsec:lym-sperner}

$\mathcal{F}\subseteq 2^{[n]}$ is an antichain $\Rightarrow$ $\sum_{F\in\mathcal{F}}1/\binom{n}{|F|}\le 1$. Consequence: $|\mathcal{F}|\le\binom{n}{\lfloor n/2\rfloor}$.

\subsection{Burnside's lemma}
\label{subsec:burnside-combinatorics}

A finite group $G$ acts on $X$ $\Rightarrow$ number of orbits $=\frac{1}{|G|}\sum_{g\in G}|X^g|$.

\subsection{Hall's theorem (on matchings)}
\label{subsec:hall}

$A_1,\dots,A_n$ have an SDR $\iff$ for all $S\subseteq[n]$: $|\bigcup_{i\in S}A_i|\ge|S|$.

\subsection{Cayley's formula and Pr\"ufer code}
\label{subsec:cayley-formula}

Number of labeled trees on $n$ vertices $=n^{n-2}$.

\medskip

\section{Linear algebra methods}
\label{sec:linalg-combinatorics}

\subsection{Dimensionality argument}
\label{subsec:dimensionality}

$v_1,\dots,v_m$ are linearly independent in $V$ of dimension $n$ $\Rightarrow$ $m\le n$.

\subsection{Diagonal independence criterion}
\label{subsec:diagonal-criterion}

If $f_i(x_j)\ne 0$ at $i=j$ and $=0$ at $i\ne j$, then $f_i$ are linearly independent.

\subsection{Oddtown}
\label{subsec:oddtown}

$A_i\subseteq[n]$, $|A_i|$ odd, $|A_i\cap A_j|$ even for $i\ne j$ $\Rightarrow$ $m\le n$.

\subsection{Eventown}
\label{subsec:eventown}

$|A_i|$ and $|A_i\cap A_j|$ even for all $i,j$ $\Rightarrow$ $m\le 2^{\lfloor n/2\rfloor}$.

\subsection{Fisher's inequality}
\label{subsec:fisher}

$A_i\subseteq[n]$, $|A_i\cap A_j|=\lambda\ge 1$ for $i\ne j$ $\Rightarrow$ $m\le n$.

\subsection{Combinatorial Nullstellensatz (Alon)}
\label{subsec:combinatorial-nullstellensatz-comb}

$f\in\mathbb{F}[x_1,\dots,x_n]$ of degree $d=\sum t_i$, coefficient of $\prod x_i^{t_i}\ne 0$. $S_i\subseteq\mathbb{F}$, $|S_i|>t_i$ $\Rightarrow$ $\exists s\in\prod S_i$ with $f(s)\ne 0$ (see also \ref{subsec:combinatorial-nullstellensatz-poly}).

\medskip

\section{Additive combinatorics}
\label{sec:additive-combinatorics}

\subsection{Cauchy--Davenport theorem}
\label{subsec:cauchy-davenport}

For non-empty $A,B\subseteq\mathbb{Z}/p\mathbb{Z}$ ($p$ prime):
$$|A+B|\ge\min(p,\,|A|+|B|-1).$$

\textit{When to use.} Lower bounds on sumsets, existence of solutions $a+b=c$ in subsets of $\mathbb{F}_p$.

\textit{Idea.} Induction on $|A|$: shift and delete intersection.

\subsection{Ruzsa triangle inequality}
\label{subsec:ruzsa-triangle}

For non-empty $A,B,C\subseteq G$ (abelian group):
$$|A+C|\cdot|B|\le|A+B|\cdot|B+C|.$$

\textit{Corollary.} $|A-A|\le\bigl(|A+A|/|A|\bigr)^2|A|$.

\subsection{Petridis's lemma}
\label{subsec:petridis}

For $A,B\subseteq G$ with $|A+B|\le K|A|$ there exists non-empty $X\subseteq A$ with $|X+B'|\le K|X|$ for any $B'\subseteq B$, where $B'$ is an arbitrary subset.

\subsection{Pl\"unnecke--Ruzsa inequality}
\label{subsec:plunnecke-ruzsa}

If $|A+A|\le K|A|$, then $|nA-mA|\le K^{n+m}|A|$, where $nA = A+\cdots+A$ ($n$ times).

\textit{When to use.} Bounds on iterated sumsets under small doubling.

\subsection{Erd\H{o}s--Ginzburg--Ziv (EGZ)}
\label{subsec:egz-combinatorics}

Among any $2n-1$ integers there exist $n$ whose sum is divisible by $n$.

\textit{Generalization (Kemnitz).} Among $2n-1$ elements of any abelian group of order $n$ there are $n$ with sum $0$ (see also \ref{subsec:egz-poly}).

\medskip

\section{Generating functions (OGF and EGF)}
\label{sec:generating-functions}

\subsection{Basic operations}
\label{subsec:gf-basic}

OGF: $A(x)=\sum a_n x^n$: $AB$ gives convolution, $x A$ gives shift, $A/(1-x)$ gives partial sums.
EGF: $\hat A(x)=\sum a_n x^n/n!$: product gives binomial convolution.

\subsection{Roots of unity filter}
\label{subsec:roots-unity-filter}

$\sum_{n\equiv r\pmod m}a_n x^n=\frac1m\sum_{j=0}^{m-1}\omega^{-jr}A(\omega^j x)$ with $\omega=e^{2\pi i/m}$.

\subsection{Linear recurrences with constant coefficients}
\label{subsec:linear-recurrence}

$\sum c_i a_{n+i}=0$ $\Rightarrow$ $A(x)=P(x)/Q(x)$, where $Q(x)=c_k+\dots+c_0x^k=\prod(1-\alpha_j x)^{m_j}$. Then $a_n=\sum Q_j(n)\alpha_j^n$.

\subsection{Catalan numbers}
\label{subsec:catalan}

$C_{n+1}=\sum C_k C_{n-k}$, $C_0=1$ $\Rightarrow$ $C(x)=1+x C(x)^2$, $C_n=\frac1{n+1}\binom{2n}{n}$.

\subsection{Lagrange inversion formula}
\label{subsec:lagrange-inversion}

$w=x\phi(w)$, $\phi(0)\ne 0$ $\Rightarrow$ $[x^n]H(w)=\frac1n[t^{n-1}] H'(t)\phi(t)^n$.

\medskip

\section{Recurrence relations and closed forms}
\label{sec:recurrences}

\subsection{Closed form of linear recurrences}
\label{subsec:closed-form-recurrence}

$a_n=\sum c_i a_{n-i}$, roots $\lambda_i$ of multiplicity $m_i$ $\Rightarrow$ $a_n=\sum P_i(n)\lambda_i^n$, $\deg P_i\le m_i-1$. Non-homogeneous case: particular solution $n^s R(n)\mu^n$ with $s$ equal to the multiplicity of $\mu$ as a characteristic root.

\subsection{Companion matrix and powers}
\label{subsec:companion-matrix}

$\mathbf v_n=C\mathbf v_{n-1}$, $C^n$ can be computed in $O(\log n)$ multiplications.

\subsection{Derangements}
\label{subsec:derangements}

$D_n=n!\sum_{k=0}^n(-1)^k/k!$, $D_n=nD_{n-1}+(-1)^n$. EGF: $e^{-x}/(1-x)$.

\subsection{Stirling numbers of the second kind}
\label{subsec:stirling-2}

$S(n,k)=kS(n-1,k)+S(n-1,k-1)$. $S(n,k)=\frac1{k!}\sum_{j=0}^k(-1)^j\binom{k}{j}(k-j)^n$. EGF: $(e^x-1)^k/k!$.

\subsection{Fibonacci identities}
\label{subsec:fibonacci}

$F_{n-1}F_{n+1}-F_n^2=(-1)^n$, $F_{m+n}=F_mF_{n+1}+F_{m-1}F_n$. Consequence: $\gcd(F_m,F_n)=F_{\gcd(m,n)}$.

\medskip

\section{Combinatorial geometry}
\label{sec:combinatorial-geometry}

\subsection{Helly's theorem}
\label{subsec:helly}

Convex sets $K_1,\dots,K_m\subseteq\mathbb{R}^d$, $m\ge d+1$. If any $d+1$ of them have a common point, then $\bigcap_{i=1}^m K_i\neq\varnothing$.

\textit{When to use.} Problems about coverings, intersections of convex figures, point arrangements.

\subsection{Sylvester--Gallai theorem}
\label{subsec:sylvester-gallai}

For a finite set of points in the plane not all on one line, there exists a line containing exactly two of them.

\textit{Idea.} Consider the pair (point, line) with minimal distance; if the line has $\ge3$ points, another pair with smaller distance is found.

\subsection{Minkowski's convex body theorem}
\label{subsec:minkowski-convex}

$K\subset\mathbb{R}^n$ is a convex centrally symmetric ($K=-K$) body with $\operatorname{vol}(K)>2^n$. Then $K$ contains a non-zero integer point $z\in\mathbb{Z}^n\setminus\{0\}$.

\textit{When to use.} Diophantine approximation, existence of solutions to systems of inequalities, quadratic forms.

\medskip

\section{Inclusion--exclusion and M\"obius inversion}
\label{sec:inclusion-exclusion}

\subsection{Inclusion--exclusion principle (PIE)}
\label{subsec:pie}

$|\bigcap\overline{A_i}|=\sum_{k=0}^n(-1)^k S_k$, where $S_k=\sum_{|I|=k}|\bigcap_{i\in I}A_i|$. Works for measures/probabilities.

\subsection{Bonferroni inequalities}
\label{subsec:bonferroni}

Even truncation $m=2r$: $|\bigcup A_i|\ge\sum_{k=1}^{2r}(-1)^{k+1}S_k$. Odd $m=2r+1$: $|\bigcup A_i|\le\sum_{k=1}^{2r+1}(-1)^{k+1}S_k$.

\subsection{M\"obius inversion}
\label{subsec:mobius-inversion}

On a locally finite poset: $g(y)=\sum_{x\le y}f(x)$ $\iff$ $f(y)=\sum_{x\le y}\mu(x,y)g(x)$. For the Boolean lattice: $\mu(I,J)=(-1)^{|J\setminus I|}$. For the divisor lattice: $\mu(d,n)=\mu(n/d)$.

\subsection{Euler's totient $\varphi$}
\label{subsec:euler-totient}

$\varphi(n)=n\prod_{p\mid n}(1-1/p)=\sum_{d\mid n}\mu(d)n/d$. M\"obius inversion: $\sum_{d\mid n}\varphi(d)=n$.

\subsection{Chromatic polynomial}
\label{subsec:chromatic}

$P(G,x)=\sum_{S\subseteq E}(-1)^{|S|}x^{c(V,S)}$ (inclusion–exclusion), where $c(V,S)$ is components of $(V,S)$. Recurrence: $P(G,x)=P(G\setminus e,x)-P(G/e,x)$.
